\documentclass[10pt]{article}
\usepackage{amsmath, amssymb, amsthm}
\usepackage{geometry}
\usepackage{hyperref}
\usepackage{booktabs}
\usepackage{setspace}
\usepackage{titlesec}
\usepackage{enumitem}
\usepackage{fontspec}
\usepackage{tikz}
\usetikzlibrary{calc}
\usepackage{eso-pic}
\usepackage[useregional]{datetime2}

\newcommand{\mymarginlabel}{DRAFT: David Ewing dew59@uclive.ac.nz +64276427681 | Compiled: \DTMnow}

\AddToShipoutPictureBG{%
  \begin{tikzpicture}[remember picture,overlay]
    \node[rotate=90, anchor=north]
      at ($(current page.north west)+(1.4cm,-13cm)$) {\scriptsize\ttfamily \mymarginlabel};
  \end{tikzpicture}%
}

\geometry{a4paper,top=2cm,bottom=2cm, left=3.2cm,right=3.2cm}
\setlength{\parindent}{0pt}
\setlength{\parskip}{0.8em}
\setstretch{1.15}
\setmainfont{Libertinus Serif}

\titleformat{\section}{\large\bfseries}{\thesection}{0.5em}{}
\titleformat{\subsection}{\normalsize\bfseries}{\thesubsection}{0.5em}{}

\hypersetup{
    colorlinks=true,
    linkcolor=black,
    urlcolor=blue,
}

\title{Proposal: Unified Token Architecture for Large-Scale CVR Prediction\\
\large QQ Algorithm Challenge --- TAAC \& KDD 2026\\
\normalsize \href{https://algo.qq.com/\#intro}{algo.qq.com}}
\author{Ben Tsai \and Yu Xia \and David Ewing} 
\date{April 2026 \\ \small Proposal v0.01 --- For Team Review}

\begin{document}
\maketitle

\begin{abstract}
I am writing this proposal to outline how we might approach the QQ
Algorithm Challenge. These are just a set of ideas I am suggesting,
rather than starting with a blank sheet. I invite everyone to review,
challenge, and/or agree on the direction before we begin
implementation --- but whatever we decide, let us keep it in a current
document, even if it changes, until such time as we agree we understand
it well enough that we no longer need to keep a living document updated.

The contest goal appears to be providing a unification of Sequence
Modelling and Feature Interaction for Large-scale Recommendation. It
appears to be a reaction to the fact that two concepts have been
diverging over time, and that there may be some benefit in putting
something together that unifies them:
\begin{itemize}
    \item \textbf{Feature-interaction models} --- high-dimensional
          multi-field categorical and contextual features
          (FM\footnote{FM: Factorisation Machine. Models pairwise
          feature interactions via inner products of learned embedding
          vectors. Rendle, 2010.},
          DeepFM\footnote{DeepFM: Deep Factorisation Machine. Combines
          an FM component with a deep neural network to capture both
          low-order and high-order feature interactions. Guo et al.,
          2017.},
          DCN\footnote{DCN: Deep \& Cross Network. Uses explicit cross
          layers to learn bounded-degree feature interactions efficiently
          alongside a deep network. Wang et al., 2017.},
          DIN\footnote{DIN: Deep Interest Network. Uses an attention
          mechanism to adaptively weight historical user behaviours
          relative to the target item. Zhou et al., 2018.}).
    \item \textbf{Sequential models} --- temporal dynamics of user
          behaviour via Transformer-style ranking models
          (SASRec\footnote{SASRec: Self-Attentive Sequential
          Recommendation. Applies a unidirectional Transformer to model
          sequential user behaviour for next-item prediction. Kang \&
          McAuley, 2018.},
          BERT4Rec\footnote{BERT4Rec: Sequential Recommendation with
          Bidirectional Encoder Representations from Transformers.
          Adapts the BERT masked-language-model objective to learn
          bidirectional user behaviour representations. Sun et al.,
          2019.}).
\end{itemize}
I propose a unified token representation that places sequential
behaviour events and non-sequential multi-field features into a single
homogeneous token stream, processed by a shared stackable Transformer
backbone. I believe this approach directly addresses both the
leaderboard metric (AUC of ROC on CVR prediction) and the two
innovation awards (Unified Block Innovation; Scaling Law Innovation).
\end{abstract}

\noindent\textbf{Keywords:} unified architecture, tokenisation, CVR prediction,
transformers, feature interaction, scaling laws, recommender systems

\newpage

\section{Problem Framing}

As I see it from a bit of research and from the goals of this contest, industrial recommender systems have evolved into two 
independent paradigms:

\begin{itemize}[noitemsep]
    \item \textbf{Feature-interaction models} (FM, DeepFM, DCN, DIN) ---
          focus on high-dimensional multi-field categorical and contextual
          features; optimise cross-feature interactions.
    \item \textbf{Sequential models} (GRU4Rec, SASRec, BERT4Rec,
          Transformer-based rankers) --- focus on temporal dynamics of user
          behaviour; encode ordered item sequences.
\end{itemize}

Both paradigms have achieved  success, yet they
have been optimised separately. In my view, this creates three
structural problems that I believe we should address:

\begin{enumerate}[noitemsep]
    \item \textbf{Shallow cross-paradigm interaction.} Concatenation of a
          sequence embedding with feature embeddings discards the
          cross-attention signal between the two token types.
    \item \textbf{Inconsistent objectives.} Joint training of two
          architecturally distinct sub-networks introduces competing
          gradient signals and higher engineering complexity.
    \item \textbf{Limited scalability.} Separate encoders do not scale
          uniformly; hardware efficiency is compromised by heterogeneous
          compute paths.
\end{enumerate}
%\newpage
\section{Proposed Approach: UniToken Architecture}

\subsection{Core Idea}

The central idea is simple: stop treating user history and user
attributes as two separate inputs that get glued together at the end.
Instead, represent everything --- past interactions, profile fields, and
the candidate item --- as items in a single list, and let one shared
network read the whole list at once.

Concretely, I propose assembling a single sequence of tokens
$\mathcal{T}$:
\begin{itemize}[noitemsep]
    \item the user's recent interactions, ordered by time:
          $[t_1, t_2, \ldots, t_L]$,
    \item the user's static feature fields (age group, region, device,
          etc.): $[f_1, f_2, \ldots, f_K]$,
    \item the candidate item being scored: $[q]$.
\end{itemize}
All three groups are placed end-to-end into one list and fed into a
single Transformer. The $\|$ symbol below just means
``join these lists one after the other'':

\[
\mathcal{T} = \underbrace{[t_1, t_2, \ldots, t_L]}_{\text{behaviour sequence}}
\;\|\;
\underbrace{[f_1, f_2, \ldots, f_K]}_{\text{non-sequential feature fields}}
\;\|\;
\underbrace{[\,q\,]}_{\text{target item / query}}
\]

Every token in the list has the same size (embedding dimension $d$), so
the Transformer can treat them all uniformly.

\textbf{Is this a time series approach?} No --- and the distinction is
important for the team. A time series model (e.g.\ ARIMA or an LSTM
forecaster) asks ``how will this measurement change over time?'' and
tries to predict the next value of the same signal. What we are doing is
different: the ordered history $[t_1, \ldots, t_L]$ provides context
about what the user has found relevant in the past, but the question we
are answering is ``will this user convert on this specific item $q$?'' ---
a binary yes/no classification, not a forecast. The Transformer does not
model whether behaviour is trending up or down; it models which parts of
the history are \emph{relevant} to the candidate item. This distinction
affects how we set up positional embeddings and the training loss.
\newpage
\subsection{Unified Token Embedding}

Before placing everything into one list, each token needs to be
converted into a fixed-length vector of numbers --- this is called an
\emph{embedding}. The key point is that every token, regardless of
whether it is a past interaction, a feature field, or the candidate
item, is converted using the \emph{same} process. Nothing is
special-cased.

Each token's vector is the sum of three parts:
\begin{enumerate}[noitemsep]
    \item \textbf{What it is} --- a learned vector for the token's
          identity (e.g.\ ``item~42'', ``region: NZ'').
    \item \textbf{Where it sits} --- a learned vector encoding its
          position: time order for behaviour tokens, field index for
          feature tokens.
    \item \textbf{What type it is} --- a learned vector marking
          whether this token is from the behaviour sequence, the
          feature fields, or is the query item.
\end{enumerate}
Summing these three parts gives the final vector $\mathbf{e}_i$ for
token $i$:

\[
\mathbf{e}_i = \mathrm{Embed}(v_i) + \mathbf{p}_i + \mathbf{s}_i
\]

All three parts have the same size $d$, so adding them is
straightforward. The type marker $\mathbf{s}_i$ is what lets the
Transformer tell the three groups apart, even though they share the
same architecture.

\subsection{Shared Stackable Backbone}

Once every token has been converted to a vector, the full list is passed
through a stack of $N$ identical Transformer blocks. Each block does two
things:
\begin{enumerate}[noitemsep]
    \item \textbf{Attention} --- every token looks at every other token
          and decides how much to ``borrow'' information from it. This
          is how a past interaction can influence how the model reads a
          feature field, and vice versa.
    \item \textbf{Feed-forward update} --- each token's vector is
          updated independently using a small two-layer network.
\end{enumerate}
Critically, \emph{every block treats all token types identically}. There
is no separate block for behaviour tokens and another for feature tokens
--- just $N$ identical blocks stacked on top of each other, each
attending across the entire list:

\[
\mathbf{H}^{(n)} = \mathrm{TransformerBlock}\!\left(\mathbf{H}^{(n-1)}\right),
\qquad n = 1,\ldots,N
\]

$N$ is a design choice we can tune --- more blocks means more capacity
but also more compute. The scaling experiments in
Section~\ref{sec:eval} are designed to understand this trade-off.

\textbf{Implementation note --- non-reentrant blocks.} Each of the $N$
blocks will be a \emph{separate} instance with its own independently
learned weights, not a single block whose weights are reused across
passes. This is the standard Transformer design (as in BERT and GPT)
and the reason it matters here is twofold. First, separate weights
allow each layer to learn qualitatively different transformations of
the token representations --- early layers may resolve local token
relationships, later layers higher-order ones. Second, and more
practically for this contest, the depth-scaling ablations
($N = 1, 2, 4, 8$) only produce a meaningful parameter-count and
compute curve if each depth level genuinely adds new parameters. If a
single block's weights were shared across all $N$ passes (a design
sometimes called a Universal Transformer), then $N = 8$ and $N = 1$
would have identical parameter counts, and any scaling law we reported
would be an artefact of inference compute rather than model capacity ---
which is not the claim we want to make for the Scaling Law Innovation
Award.
\newpage
\subsection{Prediction Head}

After the $N$ Transformer blocks have processed the full token list, we
only need to look at one token's output vector: the query token $q$
(the candidate item). By this point, that vector has been enriched by
attending to the user's entire history and all feature fields.

We pass it through a single linear layer followed by a sigmoid function,
which squashes the result to a number between 0 and 1 --- our predicted
conversion probability:

\[
\hat{y} = \sigma\!\left(\mathbf{w}^\top \mathbf{h}_q^{(N)} + b\right).
\]

In plain terms: if $\hat{y}$ is close to~1, the model thinks the user
will convert; close to~0, it thinks they will not. We train by comparing
$\hat{y}$ to the actual outcome (1~=~converted, 0~=~did not) using
binary cross-entropy loss.

\subsection{Where This Sits Relative to Prior Work}

I cannot yet claim with confidence that this
architecture is the first of its kind. What I can say is that ChatGPT has shown me that: 

\begin{itemize}[noitemsep]
    \item DIN, SIM, TWIN --- use the target item to gate the behaviour
          sequence via attention, but still process feature fields
          through a separate encoder.
    \item ETA, HSTU --- introduce structural changes to the sequence
          encoder (e.g.\ efficient attention, hierarchical units) but,
          as far as I can determine, do not place feature fields into
          the same token stream as the sequence.
\end{itemize}

The contest description also references papers [1--3] as work that
``begins to bridge the two branches''. I have not yet read these in
detail. \textbf{Before we can make any novelty claim, we should read
those papers and check whether any of them already proposes a fully
unified tokenisation.}

\subsubsection*{Placeholder: Literature Comparison Table}

\textit{This subsection is a placeholder. Once the team has reviewed
the contest references [1--3] and any other closely related work, we
should replace this with a short comparison table along the lines of:}

\medskip
\begin{tabular}{@{}llll@{}}
\toprule
Approach & Unified token stream? & Shared backbone? & Notes \\
\midrule
DIN / SIM / TWIN & No & No & Separate seq.\ encoder \\
ETA / HSTU       & Partial & No & Seq.\ only \\
{[}1{]}--{[}3{]} & TBD & TBD & To be reviewed \\
\textbf{UniToken (ours)} & \textbf{Yes} & \textbf{Yes} & \textit{claim to verify} \\
\bottomrule
\end{tabular}
\medskip

\textit{If any of [1--3] already proposes the same scheme, the novelty
claim in this proposal will need to be revised accordingly.}

\section{Evaluation Plan}
\label{sec:eval}

\subsection{Primary Metric}
AUC of ROC on the competition CVR prediction task.
\newpage
\subsection{Ablations}
Table~\ref{tab:ablations} describes the planned ablation study to isolate
the contribution of each design choice.

\begin{table}[h]
\centering
\caption{Planned ablations}
\label{tab:ablations}
\begin{tabular}{@{}lll@{}}
\toprule
Variant & Change from full model & Expected effect \\
\midrule
Separate encoders & Split backbone per token type & AUC drop \\
No segment embeddings & Remove $\mathbf{s}_i$ & Minor AUC drop \\
No cross-type attention & Mask across token types & Large AUC drop \\
Depth scaling ($N=1,2,4,8$) & Vary backbone depth & Scaling law curve \\
Width scaling ($d=32,64,128,256$) & Vary embedding dim & Scaling law curve \\
\bottomrule
\end{tabular}
\end{table}

\subsection{Scaling Law Exploration}
I intend the depth and width ablations above to produce a
loss-vs-compute curve. If a power-law relationship is observed, I
believe this directly supports a submission to the \textbf{Scaling Law
Innovation Award}, and I would like us to agree on who leads this work.

\section{Innovation Award Targets}

\begin{description}[noitemsep]
    \item[Unified Block Innovation Award (\$45,000)] I believe the
          unified token stream with a single shared backbone constitutes
          the primary novel claim. The novelty rests on the tokenisation
          scheme and the absence of type-specific modules.
    \item[Scaling Law Innovation Award (\$45,000)] I propose that we
          undertake systematic depth and width scaling experiments,
          combined with fixed compute budget analysis, as our scaling law
          contribution. These awards are independent of leaderboard rank,
          which means we can target them regardless of AUC standing.
\end{description}
%\newpage
\section{Implementation Plan}

I propose to the team the following sequence of steps, subject to team agreement:

\begin{enumerate}[noitemsep]
    \item \textbf{Baseline} --- We reproduce a standard two-tower
          (sequence + feature) model on the competition data to establish
          an AUC baseline we can all verify.
    \item \textbf{UniToken v0.1} --- We  implement the unified
          tokenisation and shared backbone (TensorFlow/Keras, building on
          \texttt{prototype-v0.02}), and share for team review.
    \item \textbf{Ablations} --- we run Table~\ref{tab:ablations}
          variants together and compare results.
    \item \textbf{Scaling runs} --- we vary $N$ and $d$; log AUC and
          FLOPs to build the scaling law curve.
    \item \textbf{Submission} --- leaderboard entry and workshop paper
          draft, with contributions from all team members.
\end{enumerate}
\newpage
\section{Open Questions --- }

\begin{itemize}[noitemsep]
    \item Should we discretise continuous features (e.g.\ dwell time,
          price) into tokens, or inject them as additive scalings? I do
          not yet have a strong preference.
    \item What is the maximum sequence length $L$ the competition data
          and latency constraints allow? I would like us to establish
          this before implementation begins.
    \item Can we agree on TensorFlow/Keras as the implementation
          framework? 
    \item Who will lead the scaling law experiments, and who will lead
          the architecture implementation? I would like us to decide
          this explicitly.
\end{itemize}

\end{document}
